\documentclass[12pt]{article}

\usepackage[margin=2cm]{geometry}
\usepackage{siunitx}
\usepackage{booktabs}
\usepackage{hyperref}
\usepackage{xcolor}
\usepackage{amsmath}
\usepackage{graphicx}

% Define the color used for section headers
\definecolor{primary}{RGB}{0, 75, 150}

\begin{document}


\section{Performance Benchmarks}

\begin{center}
    {\large\bfseries\color{primary} Performance Comparison Report} \\[0.2cm]
    {\subsection*{Advanced Framework Benchmarks across CPU and GPU Architectures}}
    \vspace{0.1cm}
\end{center}


\subsection{LAMMPS vs. MCCB-gpu (Large-Scale Scales)}
\label{sec:large_scale}
Benchmarks were conducted using two distinct NVIDIA hardware architectures: the \textbf{RTX 2000} and the \textbf{RTX 4500}.

\begin{table}[htbp]
\centering
\caption{Simulation Runtime (Seconds) by Particle Count ($N$) for Large Scales}
\label{tab:lammps_vs_gpumc}
\small
\begin{tabular}{
    @{}
    l 
    S[table-format=2.1] 
    S[table-format=2.1] 
    @{\hspace{0.5cm}} 
    S[table-format=3.1] 
    S[table-format=2.1] 
    @{}
}
\toprule
& \multicolumn{2}{c}{\textbf{LAMMPS Time}} & \multicolumn{2}{c}{\textbf{MCCB-gpu Time}} \\
\cmidrule(r{0.5cm}){2-3} \cmidrule(l){4-5}
\textbf{Particle Count ($N$)} & {\textbf{RTX 2000}} & {\textbf{RTX 4500}} & {\textbf{RTX 2000}} & {\textbf{RTX 4500}} \\
\midrule
\num{492450}  &  7.0 &  5.0 &  18.0 &  3.8 \\
\num{1969920} & 22.0 & 18.0 &  66.0 & 13.2 \\
\num{3939840} & 44.0 & 34.0 & 123.0 & 25.6 \\
\bottomrule
\end{tabular}
\end{table}

\subsection{gpMC (CPU) vs. MCCB-gpu (GPU) Scaling (\texorpdfstring{$NVT$ and $NpT$}{NVT and NpT} Ensembles)}
\label{sec:small_scale}
This benchmark assesses execution runtime across both canonical ($NVT$) and isobaric-isothermal ($NpT$) ensembles, scaling from \num{8000} to \num{125000} particles for a binary Lennard-Jones mixture ($N_1 : N_2 = 2 : 1$, $T^* = 3.0$, $r_c = 2.5\sigma$, with isotropic volume trial moves executed every cycle in $NpT$). We compare the CPU-driven implementation (\textbf{gpMC}, compiled with Intel Fortran \texttt{ifx} 2025.2.0) executed on an \textbf{Intel(R) Xeon(R) Gold 6548Y+} (32 cores / 64 threads, 2.20 GHz) against GPU acceleration (\textbf{MCCB-gpu} Version 2.8.0) on an \textbf{NVIDIA RTX PRO 4500 Blackwell} (32 GB VRAM) for workloads of 100 Monte Carlo sweeps.

\paragraph{Trial Displacements per Particle Normalization.}
A critical distinction between the two codes lies in the operational definition of a Monte Carlo cycle:
\begin{itemize}
    \item In \textbf{gpMC}, each cycle attempts exactly one single-particle displacement per particle ($\mu_{\text{CPU}} = 1.0$). Over 100 cycles, exactly 100 trial moves per particle are evaluated ($100 \times N$ total moves).
    \item In \textbf{MCCB-gpu}, the GPU checkerboard decomposition attempts $N_{\text{move}} = 125$ trial displacements per active cell across the 8 sublattices during each cycle, yielding a total of $N_{\text{trans}} = N_{\text{cells}} \times N_{\text{move}}$ moves per cycle. The number of trial displacements attempted \emph{per particle per cycle} is therefore:
    \begin{equation}
        \mu_{\text{GPU}} = \frac{N_{\text{cells}} \times N_{\text{move}}}{N}
    \end{equation}
    yielding $\mu_{\text{GPU}} = 3.375$, $1.728$, $1.953$, and $2.744$ moves/particle/cycle for $N=\num{8000}$, \num{15625}, \num{32768}, and \num{125000}, respectively.
\end{itemize}
Consequently, over 100 simulation sweeps, \texttt{MCCB-gpu} evaluates between $1.73\times$ and $3.38\times$ \textbf{more trial displacements per particle} than \texttt{gpMC}. To provide an unbiased comparison of actual Monte Carlo sampling throughput, Table~\ref{tab:gpmc_vs_mccbgpu} reports both the raw measured wallclock time $t_{\text{raw}}$ and the normalized GPU time $t_{\text{norm}} = t_{\text{raw}} / \mu_{\text{GPU}}$ (scaled to an equivalent workload of 100 moves per particle), alongside the corresponding effective speedup:
\begin{equation}
    S_{\text{effective}} = \frac{t_{\text{CPU}}}{t_{\text{norm, GPU}}} = S_{\text{raw}} \times \mu_{\text{GPU}}
\end{equation}

\begin{table}[htbp]
\centering
\caption{Simulation Runtime (Seconds) and Speedup across CPU and GPU Architectures ($NVT$ and $NpT$ Ensembles), Accounting for Differing Trial Displacements per Particle}
\label{tab:gpmc_vs_mccbgpu}
\resizebox{\textwidth}{!}{%
\begin{tabular}{
    @{}
    l 
    c 
    @{\hspace{0.4cm}}
    c 
    c 
    @{\hspace{0.4cm}}
    c 
    c 
    @{\hspace{0.4cm}}
    c 
    c 
    @{}
}
\toprule
& \textbf{Moves/Part.} & \multicolumn{2}{c}{\textbf{gpMC CPU Time (s)}} & \multicolumn{2}{c}{\textbf{MCCB-gpu GPU Time (s)}} & \multicolumn{2}{c}{\textbf{Effective Speedup}} \\
\cmidrule(lr){3-4} \cmidrule(lr){5-6} \cmidrule(l){7-8}
\textbf{Particles ($N$)} & \textbf{($\mu_{\text{GPU}}$ / $\mu_{\text{CPU}}$)} & {\textbf{$NVT$}} & {\textbf{$NpT$}} & {\textbf{$NVT$ (Norm)$^\dagger$}} & {\textbf{$NpT$ (Norm)}} & {\textbf{$NVT$}} & {\textbf{$NpT$}} \\
\midrule
\num{8000}   & 3.375 / 1.0 &   9.95 &  15.16 & 4.97 (1.47) &  6.26 (1.85) & \textbf{6.76$\times$} ($2.00\times$) & \textbf{8.17$\times$} ($2.42\times$) \\
\num{15625}  & 1.728 / 1.0 &  23.59 &  35.95 & 8.73 (5.05) & 11.88 (6.88) & \textbf{4.67$\times$} ($2.70\times$) & \textbf{5.23$\times$} ($3.03\times$) \\
\num{32768}  & 1.953 / 1.0 &  38.63 &  59.23 & 8.05 (4.12) &  9.94 (5.09) & \textbf{9.37$\times$} ($4.80\times$) & \textbf{11.64$\times$} ($5.96\times$) \\
\num{125000} & 2.744 / 1.0 & 175.91 & 254.09 & 7.86 (2.86) & 10.14 (3.70) & \textbf{61.41$\times$} ($22.39\times$) & \textbf{68.76$\times$} ($25.06\times$) \\
\bottomrule
\addlinespace[0.2cm]
\multicolumn{8}{p{17cm}}{\footnotesize Measured on node \texttt{ladon28} (Intel Xeon Gold 6548Y+ vs NVIDIA RTX PRO 4500 Blackwell) for 100 MC sweeps. In \texttt{gpMC}, each sweep attempts exactly $1.0$ displacement per particle ($\mu_{\text{CPU}} = 1.0$). In \texttt{MCCB-gpu}, checkerboard decomposition evaluates $N_{\text{cells}} \times N_{\text{move}}$ trial displacements per sweep, yielding $\mu_{\text{GPU}} = N_{\text{cells}} N_{\text{move}} / N$ trial moves per particle ($3.375$, $1.728$, $1.953$, and $2.744$ moves/particle/cycle for $N=\num{8000}, \num{15625}, \num{32768}, \num{125000}$, respectively). The normalized GPU runtime $t_{\text{norm}} = t_{\text{raw}} / \mu_{\text{GPU}}$ (in parentheses) reflects the execution time scaled to an equivalent $1.0$ move/particle workload (100 moves per particle total). Effective speedup is defined as $S_{\text{eff}} = t_{\text{CPU}} / t_{\text{norm}} = S_{\text{raw}} \times \mu_{\text{GPU}}$, with raw unadjusted cycle speedup in parentheses. $^\dagger$ Pure GPU kernel loop times (excluding I/O) on RTX 4500 are \SI{0.85}{\second} ($N=\num{8000}$), \SI{3.20}{\second} ($N=\num{15625}$), \SI{2.61}{\second} ($N=\num{32768}$), and \SI{1.93}{\second} ($N=\num{125000}$); on RTX 2000: \SI{1.21}{\second}, \SI{4.27}{\second}, \SI{3.49}{\second}, and \SI{4.34}{\second}.} \\
\end{tabular}%
}
\end{table}

\subsection{Thermodynamic Consistency Benchmarks (\texorpdfstring{$NpT \leftrightarrow NVT$}{NpT <-> NVT})}
\label{sec:npt_nvt_consistency}
To verify the physical fidelity and thermodynamic equivalence of the simulation ensembles, an automated two-stage workflow was developed and executed across cluster nodes \texttt{ladon28} and \texttt{ladon29}:
\begin{enumerate}
    \item \textbf{Stage 1 ($NpT$)}: An isobaric-isothermal simulation is initialized and equilibrated at target pressure $P^*_{\text{target}}$ and temperature $T^*$, recording the mean volume $\langle V \rangle_{NpT}$ (equilibrium density $\langle \rho^* \rangle_{NpT}$) and generating an exact configuration restart checkpoint.
    \item \textbf{Stage 2 ($NVT$)}: A canonical simulation is executed from the restart configuration at fixed density $\rho^* = \langle \rho^* \rangle_{NpT}$.
    \item \textbf{Validation}: Statistical equivalence $\langle U/N \rangle_{NVT} \approx \langle U/N \rangle_{NpT}$ is verified for continuous potential fluids (Lennard-Jones and Site-Site Patchy colloids), while the virial contact extrapolation $\langle P_{\text{virial}} \rangle$ is cross-validated against the analytical Carnahan-Starling equation of state for Hard Spheres.
\end{enumerate}

Simulations comprised $N = \num{7695}$ particles per system, executed concurrently under SLURM on dedicated RTX PRO 4500 GPUs.

\begin{table}[htbp]
\centering
\caption{Thermodynamic Consistency Validation across Interaction Models}
\label{tab:thermo_consistency}
\resizebox{\textwidth}{!}{%
\begin{tabular}{
    @{}
    l 
    l 
    c 
    c 
    c 
    c 
    @{}
}
\toprule
\textbf{Model} & \textbf{Observable} & \textbf{$NpT$ Ensemble} & \textbf{$NVT$ Ensemble} & \textbf{Discrepancy} & \textbf{Agreement} \\
\midrule
\textbf{Lennard-Jones} & Density $\langle \rho^* \rangle$ & $0.47230 \pm 0.00405$ & $0.47029$ (fixed) & 0.43\% & $0.50\sigma$ \\
($T^* = 1.50, P^* = 0.60$) & Energy $\langle U/N \rangle$ & $-2.64028 \pm 0.02328$ & $-2.63133 \pm 0.01100$ & \textbf{0.34\%} & \textbf{0.35$\sigma$ ($Z$-score)} \\
\midrule
\textbf{Hard Spheres} & Virial $P^*_{\text{NpT}}$ & $0.67674 \pm 0.00218$ & -- & \textbf{0.03\%} & vs. $P_{\text{CS}} = 0.67697$ \\
($T^* = 1.00$) & Virial $P^*_{\text{NVT}}$ & -- & $0.72284 \pm 0.00059$ & \textbf{0.04\%} & vs. $P_{\text{CS}} = 0.72311$ \\
\midrule
\textbf{Site-Site Patchy} & Density $\langle \rho^* \rangle$ & $0.71579 \pm 0.00627$ & $0.72375$ (fixed) & 1.11\% & Equilibrated \\
($T^* = 0.154, P^* = 0.05$) & Energy $\langle U/N \rangle$ & $-1.87301 \pm 0.00981$ & $-1.88494 \pm 0.00325$ & \textbf{0.64\%} & \textbf{1.15$\sigma$ ($Z$-score)} \\
& Restart Point $U/N$ & $-1.88562$ & $-1.88560 \pm 0.00034$ & \textbf{0.001\%} & Exact Match \\
\bottomrule
\addlinespace[0.2cm]
\multicolumn{6}{p{17cm}}{\footnotesize Evaluated for Hard Spheres (HS), Lennard-Jones (LJ), and Site-Site Patchy (SSP, 4-patch tetrahedral Palaia model) fluids with $N = \num{7695}$ particles. Virial pressure for HS is calculated via pair-distribution contact extrapolation $g(\sigma^+)$ and compared against the Carnahan-Starling equation of state $P_{\text{CS}}(\rho)$.} \\
\end{tabular}%
}
\end{table}

\section{Discussion and Structural Insights}
An analysis of the benchmarking profiles yields three key architectural and physical insights:
\begin{itemize}
    \item \textbf{Heterogeneous Paradigm Shifts and Sampling Throughput:} Transitioning from multicore CPU architectures (\texttt{gpMC}) to modern GPU acceleration (\texttt{MCCB-gpu} on RTX PRO 4500) delivers substantial speedups. On a raw cycle-by-cycle basis, GPU speedups reach \textbf{22.4$\times$} ($NVT$) and \textbf{25.1$\times$} ($NpT$) at $N = \num{125000}$. When normalizing for the actual number of trial displacements attempted per particle (since \texttt{MCCB-gpu} performs up to $3.38\times$ more trial moves per particle per cycle), the true sampling throughput speedup reaches \textbf{61.4$\times$} in the canonical ($NVT$) ensemble and \textbf{68.8$\times$} in the isobaric-isothermal ($NpT$) ensemble.
    \item \textbf{Anomalous Scaling Profile across Ensembles:} In both the $NVT$ and $NpT$ datasets, scaling from $N = \num{15625}$ to $N = \num{125000}$ exhibits an execution time reduction/plateau on the GPU (e.g., dropping from \SI{8.73}{\second} to \SI{7.86}{\second} in $NVT$, and plateauing at $\sim \SI{10}{\second}$ in $NpT$, despite an $8\times$ particle increase). This demonstrates a hardware architectural threshold where GPU block-level memory coalescing and thread warp occupancy reach peak efficiency.
    \item \textbf{Thermodynamic Rigor and Conservation:} The two-stage consistency test suite confirms that \texttt{MCCB-gpu} preserves strict thermodynamic equivalence across ensembles: continuous models (LJ and SSP) reproduce canonical energies within sub-percent statistical uncertainty ($0.34\%$ for LJ, $0.64\%$ for SSP, $Z < 1.2\sigma$), while athermal contact virial evaluation reproduces the analytical Carnahan-Starling equation of state within $\mathbf{< 0.04\%}$ accuracy.
\end{itemize}

\end{document}

